\chapter{Blood Pressure Measurement}

\section*{What Is This Test?}
Blood pressure measurement estimates the pressure exerted by circulating blood against the walls of the arteries. It reports systolic blood pressure (the pressure during ventricular contraction) over diastolic blood pressure (the pressure during ventricular relaxation), in millimetres of mercury (mmHg). Blood pressure is used to screen for hypertension, assess cardiovascular risk, and monitor treatment.

\section*{Alternative Names}
\begin{itemize}
  \item Blood Pressure Test
  \item BP Measurement
  \item Sphygmomanometry
  \item Office Blood Pressure
  \item Home Blood Pressure Monitoring
  \item Ambulatory Blood Pressure Monitoring
\end{itemize}
\section*{How It Is Performed}
The patient should rest quietly for at least 5 minutes, with the back supported, feet on the floor, legs uncrossed, and the arm supported at heart level. A validated, correctly sized upper-arm cuff is placed on bare skin. At least two readings are usually obtained one minute apart and interpreted with the clinical context. Home monitoring uses repeated readings over several days. Ambulatory monitoring records readings at programmed intervals during usual daytime and nighttime activities.

\section*{Why It Is Ordered}
\begin{itemize}
  \item Screening for elevated blood pressure and hypertension
  \item Confirming an elevated office reading with home or ambulatory monitoring when appropriate
  \item Assessing treatment response and long-term blood-pressure control
  \item Detecting white-coat hypertension, masked hypertension, and abnormal nighttime patterns
  \item Evaluating symptoms or complications associated with very high or very low blood pressure
\end{itemize}

\section*{Preparation, Risks, and Limitations}
The measurement is noninvasive. Caffeine, nicotine, exercise, pain, anxiety, a full bladder, talking, unsupported posture, and an incorrectly sized cuff can distort the result. A single reading does not establish a diagnosis in most people; repeated standardized measurements and appropriate clinical assessment are required. Markedly high readings accompanied by chest pain, severe shortness of breath, neurologic symptoms, confusion, or visual loss require urgent medical evaluation.

\section*{Interpretation}
Interpretation depends on the guideline, age, pregnancy status, comorbidities, and measurement setting. The systolic and diastolic values should be considered separately, and the higher risk category generally guides classification. Home and ambulatory readings often use different thresholds from office readings.
\section*{Understanding Results and Follow-up}
The laboratory report should identify the specimen, method, units, reference interval or decision limit, and whether the result is positive, negative, abnormal, or inconclusive. Reference intervals vary with age, sex, pregnancy, specimen type, assay, and laboratory; a result outside a range is not automatically a diagnosis, and a result within a range does not always exclude disease. Discuss unexpected results with the ordering clinician, who can decide whether repeat or confirmatory testing, treatment, or referral is appropriate. Do not change medicines or delay urgent care based on an isolated result.


\section*{Regulatory Status and Authorization Date}
This chapter describes a test category, not a specific commercial product. FDA, CDSCO, EU IVDR, and MHRA approval or registration dates therefore cannot be assigned without the assay manufacturer, product name, instrument, and intended use. For laboratory-developed tests, an individual agency approval date may not exist. See the Preface for the interpretation of regulatory status.

\clearpage
